\documentclass[12pt,a4paper]{article}
\usepackage[utf8]{inputenc}
\usepackage[T1]{fontenc}
\usepackage{amsmath,amssymb}
\usepackage{graphicx}
\usepackage{booktabs}
\usepackage{array}
\usepackage{caption}
\usepackage{hyperref}
\usepackage{geometry}
\usepackage{fancyhdr}
\usepackage{setspace}
\usepackage{xcolor}
\usepackage{float}
\usepackage{framed}

\setlength{\headheight}{14pt}
\geometry{a4paper,top=2.5cm,bottom=2.5cm,left=2.5cm,right=2.5cm}
\pagestyle{fancy}
\fancyhf{}
\rhead{\small Análisis numérico de EDPs · Proyecto Final}
\lhead{\small Pista B: Problema Inverso con PINNs}
\cfoot{\thepage}
\renewcommand{\headrulewidth}{0.4pt}
\hypersetup{colorlinks=true,linkcolor=blue!70!black,citecolor=blue!70!black}
\onehalfspacing

\newcommand{\lhat}{\hat{\lambda}}
\newcommand{\ltrue}{\lambda_{\text{real}}}
\newcommand{\Lf}{\mathcal{L}_f}
\newcommand{\Lb}{\mathcal{L}_b}
\newcommand{\Li}{\mathcal{L}_{ic}}
\newcommand{\Ld}{\mathcal{L}_d}

\begin{document}

\begin{titlepage}
  \centering\vspace*{2cm}
  {\Large\textbf{Análisis Numérico de Ecuaciones Diferenciales}\par}
  \vspace{0.5cm}{\large Proyecto Final\par}\vspace{2cm}
  {\LARGE\textbf{Redes Neuronales Informadas por la Física (PINNs):}\par}
  \vspace{0.4cm}{\LARGE\textbf{Problema Inverso de Parámetros}\par}
  \vspace{0.5cm}{\large Pista B\par}\vspace{3cm}
  {\large\textit{Ecuación de difusión 1D con coeficiente desconocido}\par}
  \vspace{3cm}
  {\normalsize Basado en: Raissi, Perdikaris y Karniadakis (2019),\\
  \textit{J. Comput. Phys.} 378:686--707\par}
  \vfill{\normalsize\today\par}
\end{titlepage}

\tableofcontents\newpage

% ============================================================
\section{Introducción y Planteamiento del Problema}
% ============================================================

En numerosos fenómenos de transporte, la forma funcional de la EDP es conocida, pero
ciertos coeficientes constitutivos permanecen inaccesibles a medición directa. El
\textbf{problema inverso de parámetros} consiste en inferirlos a partir de observaciones
dispersas y ruidosas del campo solución.

Se estudia la \textbf{ecuación de difusión unidimensional}:
\begin{equation}
  \frac{\partial u}{\partial t}=\lambda\,\frac{\partial^2 u}{\partial x^2},\qquad
  (x,t)\in[0,1]^2,\label{eq:difusion}
\end{equation}
con condición inicial $u(x,0)=\sin(\pi x)$ y condiciones Dirichlet homogéneas
$u(0,t)=u(1,t)=0$. La solución analítica exacta es
\begin{equation}
  u(x,t)=\sin(\pi x)\,e^{-\lambda\pi^2 t}.\label{eq:analitica}
\end{equation}
El parámetro $\lambda>0$ es \textbf{incógnita}; el objetivo es recuperarlo a partir de
$N_d=100$ observaciones con 1\,\% de ruido gaussiano.

\begin{figure}[htbp]
  \centering
  \includegraphics[width=0.65\textwidth]{figura_01_condicion_inicial.png}
  \caption{Condición inicial del problema: $u(x,0)=\sin(\pi x)$ en el dominio $[0,1]$.
           Esta función suave de amplitud unitaria es la referencia para todos los
           experimentos presentados en este informe.}
  \label{fig:ic}
\end{figure}

\begin{figure}[htbp]
  \centering
  \includegraphics[width=0.70\textwidth]{figura_02_observaciones_ruidosas_png.png}
  \caption{Distribución de los $N_d=100$ puntos de observación sintéticos en el
           dominio $(x,t)\in[0,1]^2$, coloreados por el valor de $u$ observado. Las
           observaciones se generaron evaluando la solución analítica~\eqref{eq:analitica}
           con $\lambda=0.1$ y añadiendo ruido gaussiano del 1\,\%.}
  \label{fig:datos}
\end{figure}

% ============================================================
\section{Metodología}
% ============================================================

\subsection{Arquitectura de la red}
Se implementó una PINN como subclase de \texttt{tf.keras.Model} en TensorFlow~2.x:
3~capas ocultas de 40~neuronas, activación $\tanh$, inicialización Glorot normal. El
coeficiente $\lhat$ se declara como \texttt{tf.Variable} escalar con restricción
$\lhat\geq10^{-6}$, inicializado en $\lhat_0=0.5$.

\subsection{Función de pérdida}
\begin{equation}
  \mathcal{L}(\theta,\lhat)=w_f\Lf+\Li+\Lb+w_d\Ld,\quad w_f=1,\;w_d=10.
\end{equation}

\subsection{Diferenciación automática}
Se usaron tapes anidados (\texttt{tape1} interno y \texttt{tape2} externo) para calcular
correctamente $\partial^2\tilde{u}/\partial x^2$.

\subsection{Entrenamiento}
Adam, $\eta=10^{-3}$, 8000~épocas. $\lhat$ y $\theta$ se optimizan conjuntamente.

% ============================================================
\section{Problema Directo}
% ============================================================

\subsection{Convergencia}

La Figura~\ref{fig:conv_directo} muestra la evolución de las componentes de la función
de pérdida durante el entrenamiento del problema directo (con $\lambda$ conocido). Las
tres curvas descienden de forma consistente, lo que confirma la convergencia simultánea
del residuo PDE y de las condiciones de contorno e iniciales.

\begin{figure}[htbp]
  \centering
  \includegraphics[width=0.88\textwidth]{figura_03_convergencia_directo_png.png}
  \caption{Evolución de la función de pérdida durante el entrenamiento del problema
           directo, descompuesta en pérdida total, residuo de la PDE y términos de
           condición inicial y de contorno (IC+BC), representada en escala logarítmica.}
  \label{fig:conv_directo}
\end{figure}

\subsection{Solución reconstruida}

La Figura~\ref{fig:sol_directo} compara la solución analítica con la predicción del
PINN y muestra el mapa de error absoluto resultante. El error máximo se concentra en
la región de mayor gradiente temporal ($t\approx0$, $x\approx0.5$), manteniéndose
dentro del orden $10^{-3}$.

\begin{figure}[htbp]
  \centering
  \includegraphics[width=0.95\textwidth]{figura_09_validacion_directo_png.png}
  \caption{Validación del problema directo: solución analítica (izquierda), predicción
           del PINN (centro) y mapa de error absoluto $|u_{\mathrm{PINN}}-u_{\mathrm{anal}}|$
           (derecha) en el dominio $[0,1]^2$.}
  \label{fig:sol_directo}
\end{figure}

% ============================================================
\section{Problema Inverso}
% ============================================================

\subsection{Convergencia del parámetro}

Las Figuras~\ref{fig:conv_lambda} y~\ref{fig:loss_inv} documentan la convergencia del
problema inverso. En particular, la Figura~\ref{fig:conv_lambda} muestra que $\lhat$
parte de $0.5$ y converge al valor real $\lambda=0.1$ en las primeras 500~épocas,
permaneciendo estable durante el resto del entrenamiento.

\begin{figure}[htbp]
  \centering
  \includegraphics[width=0.88\textwidth]{figura_05_evolucion_lambda_perdida_png.png}
  \caption{Evolución conjunta del parámetro estimado $\lhat$ (eje izquierdo, escala
           lineal) y de la pérdida total (eje derecho, escala logarítmica) durante el
           entrenamiento del problema inverso. La línea roja discontinua indica el valor
           real $\ltrue=0.1$.}
  \label{fig:conv_lambda}
\end{figure}

\begin{figure}[htbp]
  \centering
  \includegraphics[width=0.75\textwidth]{figura_04_evolucion_perdida_png.png}
  \caption{Evolución de la pérdida total durante el entrenamiento del problema inverso
           en escala logarítmica. El descenso monótono general, interrumpido únicamente
           por un transitorio en torno al paso~8, refleja la dinámica no convex del
           paisaje de optimización.}
  \label{fig:loss_inv}
\end{figure}

La Figura~\ref{fig:evol_lambda} complementa la anterior mostrando exclusivamente la
trayectoria de $\lhat$ en escala lineal, lo que permite apreciar la precisión de la
estimación final.

\begin{figure}[htbp]
  \centering
  \includegraphics[width=0.75\textwidth]{figura_06_evolucion_lambda_png.png}
  \caption{Convergencia del parámetro estimado $\lhat$ hacia el valor real
           $\lambda=0.1$ a lo largo del entrenamiento del problema inverso. El estimador
           alcanza y mantiene una proximidad inferior al 0.2\,\% a partir de las
           primeras 1000~épocas.}
  \label{fig:evol_lambda}
\end{figure}

\subsection{Resultado final}

\begin{table}[htbp]
  \centering
  \caption{Resultado del problema inverso.}\label{tab:resultado}
  \begin{tabular}{lc}\toprule
    Magnitud & Valor\\\midrule
    $\ltrue$ & 0.10000\\
    $\lhat$ estimado & 0.10017\\
    Error relativo & \textbf{0.17\,\%}\\
    $|\nabla_{\lhat}\mathcal{L}|$ final & $1.58\times10^{-6}$\\\bottomrule
  \end{tabular}
\end{table}

\subsection{Reconstrucción del campo}

La Figura~\ref{fig:sol_inverso} muestra que el campo $u(x,t)$ reconstruido con el
parámetro inferido $\lhat=0.10017$ es prácticamente indistinguible de la solución
analítica, con un error absoluto máximo inferior a $7.2\times10^{-3}$.

\begin{figure}[htbp]
  \centering
  \includegraphics[width=0.95\textwidth]{figura_10_validacion_inverso_png.png}
  \caption{Validación del problema inverso: predicción del PINN con
           $\lhat=0.10017$ (izquierda), solución analítica correspondiente (centro) y
           mapa de error absoluto $|u_{\mathrm{PINN}}-u_{\mathrm{anal}}|$ (derecha)
           en el dominio $[0,1]^2$.}
  \label{fig:sol_inverso}
\end{figure}

% ============================================================
\section{Validación}
% ============================================================

\subsection{Comparación con diferencias finitas (Crank-Nicolson)}

\begin{table}[htbp]
  \centering
  \caption{Errores $L^2$ relativos de la validación contra Crank-Nicolson.}\label{tab:cn}
  \begin{tabular}{lcc}\toprule
    Comparación & Problema Directo & Problema Inverso ($\lhat=0.10017$)\\\midrule
    CN vs.\ analítica   & $6.711\times10^{-3}$ & $6.737\times10^{-3}$\\
    PINN vs.\ analítica & $1.470\times10^{-3}$ & $1.565\times10^{-3}$\\
    CN vs.\ PINN        & $7.453\times10^{-3}$ & $7.250\times10^{-3}$\\\bottomrule
  \end{tabular}
\end{table}

Las Figuras~\ref{fig:cn_directo} y~\ref{fig:cn_inverso} presentan cortes temporales
que permiten una inspección visual directa del acuerdo entre los tres métodos. En ambos
casos, las curvas de la solución analítica, Crank-Nicolson y PINN son prácticamente
superponibles en toda la ventana temporal.

\begin{figure}[htbp]
  \centering
  \includegraphics[width=0.95\textwidth]{figura_07_cortes_temporales_directo_png.png}
  \caption{Comparación de cortes temporales en el problema directo para
           $t\in\{0,\,0.25,\,0.5,\,0.75,\,1.0\}$: solución analítica (negro sólido),
           Crank-Nicolson (azul discontinuo) y PINN (rojo punteado). Las tres curvas
           son visualmente indistinguibles en todos los instantes considerados.}
  \label{fig:cn_directo}
\end{figure}

\begin{figure}[htbp]
  \centering
  \includegraphics[width=0.95\textwidth]{figura_08_cortes_temporales_inverso_png.png}
  \caption{Comparación de cortes temporales en el problema inverso para
           $t\in\{0,\,0.25,\,0.5,\,0.75,\,1.0\}$: solución analítica (negro sólido),
           Crank-Nicolson (azul discontinuo) y PINN inverso con $\lhat=0.10017$
           (rojo punteado).}
  \label{fig:cn_inverso}
\end{figure}

\subsection{Estudio de sensibilidad}

\begin{table}[htbp]
  \centering
  \caption{Error relativo en $\lhat$ (\%) — media sobre tres semillas (rango entre paréntesis).}
  \label{tab:sensibilidad}
  \begin{tabular}{ccc}\toprule
    $N_d$ & Ruido 0\,\% & Ruido 1\,\%\\\midrule
    20  & 0.18\,\% (0.03--0.26) & 0.35\,\% (0.11--0.52)\\
    50  & 0.23\,\% (0.17--0.29) & 0.12\,\% (0.07--0.21)\\
    100 & 0.21\,\% (0.10--0.29) & 0.65\,\% (0.53--0.80)\\
    200 & 0.19\,\% (0.14--0.28) & 0.26\,\% (0.07--0.65)\\\bottomrule
  \end{tabular}
\end{table}

\begin{figure}[htbp]
  \centering
  \includegraphics[width=0.95\textwidth]{figura_11_sensibilidad_nd_ruido_png.png}
  \caption{Estudio de sensibilidad del problema inverso: error relativo en $\lhat$
           (panel izquierdo) y error $L^2$ relativo del campo reconstruido (panel
           derecho) en función del número de observaciones $N_d$, para dos niveles
           de ruido ($\sigma\in\{0\%,\,1\%\}$). La línea gris discontinua marca el
           umbral de referencia del 1\,\%. En todos los escenarios evaluados, ambas
           métricas se mantienen por debajo de dicho umbral.}
  \label{fig:sens}
\end{figure}

\subsection{Validación de plausibilidad física}
\begin{enumerate}
  \item \textbf{Positividad:} $\lhat\leftarrow\max(\lhat,10^{-6})$ tras cada paso.
  \item \textbf{Consistencia del campo:} reproduce el decaimiento exponencial analítico.
  \item \textbf{Diagnóstico de gradiente:} $|\nabla_{\lhat}\mathcal{L}|$ consistente durante todo el entrenamiento.
\end{enumerate}

% ============================================================
\section{Discusión y Limitaciones}
% ============================================================

\subsubsection*{Paisaje de optimización no convexo.}
Se detectó un mínimo local espurio ($N_d=50$, ruido 0\,\%, error $L^2=4.12\times10^{-2}$)
no reproducido en corridas posteriores, lo que motiva el uso de múltiples inicializaciones.

\subsubsection*{Ausencia de tendencia monótona con $N_d$.}
El error en $\lhat$ oscila entre 0.18\,\% y 0.23\,\% para ruido nulo, indicando que el
cuello de botella es el optimizador, no la cantidad de datos.

\subsubsection*{Variabilidad con ruido.}
Con ruido del 1\,\%, la variabilidad inter-semilla aumenta, especialmente en $N_d=100$.

\subsubsection*{Extensión a Fokker-Planck.}
Extensión natural: inferir deriva y difusión conjuntamente siguiendo Chen et al.\ (2021).

% ============================================================
\section{Conclusiones}
% ============================================================

\begin{enumerate}
  \item $\lambda$ recuperado con error relativo \textbf{0.17\,\%} (umbral: 1\,\%).
  \item Error $L^2$ del campo: $1.57\times10^{-3}$ (objetivo: $\sim10^{-3}$).
  \item PINN supera a Crank-Nicolson ($N_x=99$, $N_t=1000$) en precisión.
  \item Robusto para $N_d\geq20$ y ruido hasta 1\,\% — 24/24 configuraciones bajo umbral.
\end{enumerate}

\begin{thebibliography}{9}
\bibitem{raissi2019} Raissi, M., Perdikaris, P., \& Karniadakis, G.\,E. (2019).
Physics-informed neural networks. \textit{J. Comput. Phys.}, 378, 686--707.
\bibitem{karniadakis2021} Karniadakis, G.\,E. et al.\ (2021). Physics-informed machine learning.
\textit{Nature Reviews Physics}, 3, 422--440.
\bibitem{chen2021} Chen, X. et al.\ (2021). Solving inverse stochastic problems via
Fokker-Planck and PINNs. \textit{SIAM J. Sci. Comput.}
\bibitem{krishnapriyan2021} Krishnapriyan, A.\,S. et al.\ (2021). Failure modes in PINNs.
\textit{NeurIPS}, 34.
\end{thebibliography}

\end{document}
